% RECOVERY32: reconstructed from the immutable f3c084c0 cycle31 proof.
\statement{A.100}{Proved}{The full-GT pentagon forces the full coefficient equality}
For an actual negative involution \(\alpha\in\widehat{GT}\), the
actual regular-frame unit and canonical dihedral exponent satisfy
\[
 a_{\rm reg}=d_\alpha=2k_\alpha+1
             \quad\hbox{in }\widehat{\mathbb Z}.
 \tag{A.100.1}
\]
Equivalently the actual regular scalar is
\(\xi=-(d_\alpha-1)/4\); the full abelianization of
\(\Omega_\alpha\) is zero. This does not assert the full nonlinear
cusp equation, and retains the full-GT hypothesis.

\begin{proof}
Let u be the actual canonical norm word. A.37 proves the full
pentagon for this very word: the fixed arrow Au, its five genuine
operadic insertions, and C.28/A.36's trivial fixed subgroup in
\(\widehat P_4\) kill the entire discrepancy. Thus the literal
original marking gives
\[
 u(a,ce)u(bc,f)=u(c,f)u(ab,ef)u(a,c).
 \tag{A.100.2}
\]
Fixed arrows alone are not being substituted for that last theorem.
In the original first-cusp module
\(\Lambda=\mathbb F_2[[\widehat{\mathbb Z}]]\), with
\(x\mapsto(1,1),y\mapsto(0,T)\), put
\(U=[u(x,y)]\), \(\bar U(T)=U(T^{-1})\). The exact A.88
normalization is
\[
 \theta=1+(T-1)(U+\bar U),\quad
 (T-1)D_b=T^b-1,\qquad \theta D_{a_{\rm reg}}=1.
 \tag{A.100.3}
\]
In particular theta is a unit. Write \(k=k_\alpha\) for the
actual universal dihedral image \(u\mapsto R^k\), so
\(d_\alpha=2k+1\). The dyadic equality is already A.89.

Let \(\Bbbk\) be a finite field of characteristic two and
\(\rho\in\Bbbk^\times\) have odd order n greater than one.
Set \(t=\rho+1\), and on \(V=\Bbbk^2\) put
\[
 D=\begin{pmatrix}1&0\\t/\rho&1\end{pmatrix},\qquad
 F=\begin{pmatrix}1&t\\0&1\end{pmatrix},\qquad
 \mathcal R=FD.
 \tag{A.100.4}
\]
D,F are involutions. The matrix \(\mathcal R\) has determinant
one and distinct eigenvalues \(\rho,\rho^{-1}\), hence order n.
In \(V\rtimes\langle\rho I,D,F\rangle\), prescribe
\[
 P=(0,\rho I),\quad A=(e_1,I),\quad C=(e_2,I),
 \quad d=(e_2,D),\quad f=(e_1,F).
 \tag{A.100.5}
\]
This is a homomorphism from the entire actual A.96 source.
Indeed the three conjugation images under d are
\[
 ((\rho-1)e_2,\rho I),\quad
 (e_1+(\rho^{-1}-1)e_2,I),\quad(e_2,I),
\]
equal to \(C^{-1}PC,C^{-1}P^{-1}CPA,C\), respectively.
The three under f are
\[
 ((\rho-1)e_1,\rho I),\quad(e_1,I),
 \quad(e_2+(\rho-1)e_1,I),
\]
equal to \(A^{-1}PA,A,CPA^{-1}P^{-1}A\). Thus all six original
relations hold, and the boundary P has its nontrivial odd order.
The inverse original marking gives
\[
\begin{array}{c|c}
 a&(0,F)\\ b&(0,\rho^{-1}FD)\\ c&(0,D)\\
 d&(e_2,D)\\ e&(\rho(e_1+e_2),\rho DF)\\ f&(e_1,F).
\end{array}
\tag{A.100.6}
\]

Put \(w=e_2/t\), so \((\mathcal R-I)w=e_1\).
The pair (c,f) is the simultaneous conjugate of (D,F) by w,
whereas (a,c) is the linear pair (F,D). Their ordered dihedral
evaluations are therefore
\[
 u(c,f)=(w-\mathcal R^{-k}w,\mathcal R^{-k}),
 \qquad u(a,c)=(0,\mathcal R^k).
 \tag{A.100.7}
\]
The inverse rotation in the first formula is forced by the reversed
pair. Direct multiplication also gives
\[
\begin{gathered}
 ce=(\rho e_1+e_2,\rho F),\quad bc=(0,\rho^{-1}F),\\
 ab=(0,\rho^{-1}D),\quad ef=(e_2,\rho D).
\end{gathered}
\tag{A.100.8}
\]
For (a,ce), conjugate by the translation
\(-(e_1+e_2)/t\). Both translations become e1, with linear
heads F and \(\rho F\). On the invariant line \(\Bbbk e_1\)
the scalar pair is \((1,1),(1,\rho)\). A further scalar
translation makes the second translation zero without changing
the first. A commutator-word value has head one and its translation
is unchanged by translation conjugacy. Hence
\(u(a,ce)=(U(\rho)e_1,I)\).
For (bc,f) the scalar pair on that line is
\((0,\rho^{-1}),(1,1)\); actual duality and characteristic two
give \(u(bc,f)=(U(\rho^{-1})e_1,I)\).
The pair (ab,ef) lies over the commuting heads
\(\rho^{-1}D,\rho D\) on \(\Bbbk e_2\), so write its actual
value as \((Y_\rho e_2,I)\). This coefficient is retained, not
declared zero or identified with an unrelated parameter.

Let \(B_\rho=U(\rho)+U(\rho^{-1})\). Substitution of all five
values into the ordered pentagon yields
\[
 \mathcal R^k[I+B_\rho(\mathcal R-I)]w
                       =(1+tY_\rho)w.
 \tag{A.100.9}
\]
The vectors w and \(\mathcal Rw\) are independent. Thus w has
nonzero projections to both distinct eigenlines. Applying them to
(A.100.9) gives
\[
 \rho^k[1+(\rho-1)B_\rho]
  =1+tY_\rho
  =\rho^{-k}[1+(\rho^{-1}-1)B_\rho].
 \tag{A.100.10}
\]
Eliminating the retained Y gives precisely
\(\theta(\rho^{-1})=\rho^{2k}\theta(\rho)\), in the original
normalization (A.100.3).

The continuous geometric sums satisfy
\(\overline{D_a}=T^{1-a}D_a\). Apply bar to
\(\theta D_a=1\) and cancel its units to get
\(\bar\theta=T^{a-1}\theta\). Compare at rho with the preceding
identity and cancel the nonzero \(\theta(\rho)\). This gives
\(\rho^{a-d_\alpha}=1\). For every odd n a finite characteristic-two
field contains an element of exact order n, so
\(a\equiv d_\alpha\pmod n\) for every odd n. A.89 gives equality
on the dyadic component; together these prove (A.100.1).
Torsion-freeness of \(\widehat{\mathbb Z}\) justifies cancellation
of four in the scalar reformulation.

No four-strand canonical action or pentagon was assumed for general
GT0 data. The exact coefficient equality determines the full
abelianization of the marked defect, not its full group value.
\end{proof}
